\subsection{Further results on the axioms}\label{supp:axioms-more}

This subsection contains results on the axioms of Section~\ref{sec:axioms} that are not needed to follow the paper:
why the hypothesis of \axR{} cannot be replaced by a condition on finite values of $\LR$, why the weights of \axCM{}
cannot be given from outside, explicit formulas for the implied premium, and why monotonicity is imposed on the
cash offsets rather than on the values of the ratio. The Discussions record the author's questions that led to them.

The hypothesis of \axR{} that $X$ and $Y$ are unpriced cannot be removed, and it cannot be replaced by a condition on finite
values of $\LR$. For a loss ratio $\LR$ write $\widetilde\LR$ for the real-valued map equal to $\LR$ where $\LR$ is
finite and to $1$ elsewhere. A condition \emph{depends on $\LR$ only through its finite values} if it holds for
$\LR$ whenever it holds for some $\LR^\circ$ with $\widetilde{\LR^\circ}=\widetilde\LR$. Conditions on limits
of finite values along cash lines, such as $\lim_{s\to0}s(\LR(X+s)-1)$ or $\lim_{|s|\to\infty}s(\LR(X+s)-1)$, are of
this kind, since under \axP{} the values $\LR(X+s)$ near the singular point are finite.

\begin{proposition}[Non-finite values are needed]\label{prop:nofinite}
Let $\Pi'=\max(I_{\sqrt t},I_{\min(1,1.5t)})$ and $\LR'=\bracket{\Pi'}{\Pi'}$. Then $\LR'$ satisfies \axPe, \axP,
\axN, \axLI, \axRC{} and \axRM, its finite values are all equal to $1$, and $\LR'$ is not of the form
$\bracket{I_h}{I_g}$ with $h,g\in\Hset_1$. Consequently, no condition that depends on $\LR$ only through its finite
values and holds for every loss ratio $\bracket{I_h}{I_g}$ can replace \axR{} in Theorems~\ref{thm:A}
and~\ref{thm:B}.
\end{proposition}

\begin{proof}
$\Pi'$ is law invariant, monotone, $1$-Lipschitz and cash additive with $\Pi'(c)=c$, so the cash line of $X$ under
$\LR'$ equals $1$ except at $c=-\Pi'(X)$, where it is not finite. This gives \axPe, \axP, \axN{} and \axLI. The sure
amount $a$ with $\LR'(X-a)\notin\R$ is $\Pi'(X)$; a book $W$ never has $\LR'(W)=0$, and $\LR'(W)\notin\R$ with
$\LR'(W+1)=1$ holds exactly when $\Pi'(W)=0$. Both sure amounts in \axRC{} and \axRM{} are therefore $\Pi'(X)$, and
both axioms hold. If $\LR'=\bracket{I_h}{I_g}$, the two maps have the same non-finite books, so $I_g(X)=\Pi'(X)$ for
every $X$ by \axP, and $\Pi'$ would be comonotonic additive; Example~\ref{ex:indepR} shows that it is not. Finally,
$\widetilde{\LR'}\equiv1=\widetilde{\LR^\circ}$ for $\LR^\circ=\bracket{\E}{\E}$, which is of the form
$\bracket{I_h}{I_g}$; so every condition of the stated kind holds for $\LR'$. \qedhere
\end{proof}

On books that break even, the finite values of the cash line carry no information about the premium, and only the
location of the non-finite values does (Proposition~\ref{prop:explicit} below). The following remark shows the same
for the limit at a vanishing premium.

\begin{remark}[The limit at zero premium]\label{rem:limzero}
For $\LR=\bracket{I_h}{I_g}$, $s(\LR(X+s)-1)=s\,(I_h(X)-I_g(X))/(I_g(X)+s)$. Its limit as $s\to0$ exists for every
$X$: it is $I_h(X)$ if $X$ is unpriced and $0$ otherwise. An axiom requiring that this limit exist is therefore
automatic, and the limit is $0$ on unpriced books with $I_h(X)=0$ as well as on all books with a premium, so it
does not identify the unpriced books. Nor can \eqref{eq:R} be stated for the limits on all comonotone pairs: take
$\rho=\E$, $\Pi=I_{\sqrt t}$, $X=\1_A$ with $P(A)=\tfrac14$ and $Y=-\tfrac12$. The limits for $X$ and $Y$ are $0$,
since $\Pi(X)=\tfrac12$ and $Y$ is a sure amount, while $X+Y$ is unpriced and its limit is $\E(X+Y)=-\tfrac14$. The
hypothesis of \axR{} accounts for such pools, whose premiums cancel.
\end{remark}

\begin{discussion}[relative price]
In an earlier version the pooling property was written with the premiums inside a single fraction. The author
asked that it be written as the combination $w_{X,Y}\LR(X)+w_{Y,X}\LR(Y)$ of the two loss ratios, and that the
weight be explained as the relative price of $X$ in the comonotone portfolio $\{X,Y\}$, since this is the form in
which an underwriter reads a pooled loss ratio. The paragraph before Property~\ref{ax:CM} records that reading.
\end{discussion}

\begin{remark}[A premium given from outside]\label{rem:primitive}
Suppose that a premium functional $\Pi^\circ$ is given and \axCM{} is stated with $\Pi^\circ$ as weights. If
$\LR(X)=a\ne b=\LR(Y)$ and $\LR(X+Y)=L$, then $\Pi^\circ(X)/\Pi^\circ(Y)=(b-L)/(L-a)$, so the weights are determined
by the loss ratios. On books with $\LR=1$, however, the pooled ratio equals $1$ whatever the weights, and \axCM{}
carries no information on $\Pi^\circ$. For instance, if $\LR\equiv1$ on $\{I_g\ne0\}$, then \axCM{} holds with the
non-additive weights $\Pi^\circ=(I_g)^3$, since $a^3+b^3\ne0$ implies $a+b\ne0$. With the implied premium this
difficulty does not arise, because \axP{} locates the price point also when the ratio equals $1$ elsewhere on
the cash line.
\end{remark}

Away from the break-even set, the premium can be computed explicitly from two values of the loss ratio. Solving
\eqref{eq:pool} for $\Pi(X)$ gives the following formulas.

\begin{proposition}[Explicit formulas]\label{prop:explicit}
Assume \axP, \axN{} and \axCM, and let $X\in D$ with $\LR(X)\ne1$. Write $\phi_X(c)=\LR(X+c)$.
\begin{enumerate}[label=(\alph*),nosep]
\item \emph{Two quotes.} For every $c\notin\{0,-\Pi(X)\}$, $\phi_X(c)\ne\LR(X)$ and
  \[
   \Pi(X)=\frac{c\,\bigl(1-\phi_X(c)\bigr)}{\phi_X(c)-\LR(X)},\qquad
   \rho(X)=\LR(X)\,\frac{c\,\bigl(1-\phi_X(c)\bigr)}{\phi_X(c)-\LR(X)} .
  \]
\item \emph{Sensitivity to a sure loss.} $\phi_X$ is differentiable at $0$, $\phi_X'(0)=(1-\LR(X))/\Pi(X)\ne0$, and
  \[
   \Pi(X)=\frac{1-\LR(X)}{\phi_X'(0)},\qquad \rho(X)=\frac{\LR(X)\,\bigl(1-\LR(X)\bigr)}{\phi_X'(0)} .
  \]
\item \emph{Underwriting result.} $\Pi(X)-\rho(X)=\lim_{|c|\to\infty}c\,\bigl(1-\phi_X(c)\bigr)$.
\end{enumerate}
If $\LR(X)=1$, then $\phi_X\equiv1$ away from the price point, and $\Pi(X)$ is not determined by the finite values
of the cash line.
\end{proposition}

\begin{proof}
By Proposition~\ref{prop:readings}(a), $\phi_X(c)=(\rho+c)/(\Pi+c)$ with $\rho=\rho(X)$, $\Pi=\Pi(X)\ne0$, so
\[
 1-\phi_X(c)=\frac{\Pi-\rho}{\Pi+c},\qquad \phi_X(c)-\LR(X)=\frac{c\,(\Pi-\rho)}{\Pi\,(\Pi+c)} .
\]
Since $\LR(X)\ne1$ means $\rho\ne\Pi$, the second expression is non-zero for $c\ne0$, and the quotient gives (a).
Differentiating, $\phi_X'(c)=(\Pi-\rho)/(\Pi+c)^2$, so $\phi_X'(0)=(1-\LR(X))/\Pi$, which gives (b); and
$c\,(1-\phi_X(c))=c\,(\Pi-\rho)/(\Pi+c)\to\Pi-\rho$, which gives (c). If $\LR(X)=1$, then $\rho=\Pi$ and
$\phi_X(c)=1$ for $c\ne-\Pi$. \qedhere
\end{proof}

Formula (b) reads: the implied premium is the underwriting margin $1-\LR(X)$ divided by the sensitivity of the
loss ratio to a sure loss; formula (a) is its finite-difference version, which needs only two quoted loss ratios,
with and without a sure loss $c$. For $X=\1_A$ in Example~\ref{ex:three}, $\LR(X)=\tfrac12$ and
$\phi_X(1)=\tfrac56$, so (a) gives $\Pi(X)=\tfrac16/\tfrac13=\tfrac12$. For books outside $D$ the formulas apply to
$X+c$ with $X+c\in D$, and $\Pi(X)=\Pi(X+c)-c$ by \eqref{eq:Picash}.

The last statement of Proposition~\ref{prop:explicit} is the reason why the axioms must refer to non-finite
values of $\LR$, as \axPe{} and \axR{} do. On the break-even set the cash line carries no information about the premium except the
location of its singularity, and Remark~\ref{rem:primitive} shows that the weights in \axCM{} are then not
identified by the values of the ratio. Off the break-even set, the weights of \axCM{} could equally be written with
the explicit premium of (a) or (b) in place of $-\cinf{}$; but \axCM{} must also govern pools that contain
break-even books, and there only the price point locates the premium. Axiom \axR{} avoids weights altogether: it
moves each book to zero premium by a sure payment, which is a statement about non-finite values, and pools there.

\begin{discussion}[explicit premium formulas]
After the remark of Peng Liu reported in Section~\ref{sec:collab-revision}, the author asked whether the implicitly
defined premium could be replaced by an explicit formula in the values of $\LR$, so that the axioms could be
restated in terms of the loss ratio. Three proposals were examined. The first solves the identity \eqref{eq:pool}
for $\Pi(X)$, which gives (a) and (b) of Proposition~\ref{prop:explicit}. The second, proposed by the author, takes
the existence of $M(X):=\lim_{|c|\to\infty}c\,(1-\LR(X+c))$ as a condition and defines the premium as
$M(X)/(1-\LR(X))$; by (c) it agrees with the implied premium off the break-even set. The third, also proposed by
the author, treats a break-even book $X$ by perturbation, setting $\Pi(X):=\lim_{\varepsilon\to0}\Pi(X+\varepsilon
Y)$ for a book $Y$ such that $X+\varepsilon Y$ is outside the break-even set, with $\Pi(X+\varepsilon Y)$ given by
the second formula. The main assistant checked that each proposal agrees with the implied premium wherever it is
defined, and that the perturbation recovers it on the break-even set provided the premium is already known to be
continuous, which is part of what the axioms are meant to deliver. The conclusion was that no explicit formula can
serve as a definition on its own: on break-even books, and for $h=g$ on every book, the finite values of the cash
line carry no information about the premium. This led to the formulation of Section~\ref{sec:R}, in which the
premium enters only through the non-finite values of $\LR$.
\end{discussion}

\begin{discussion}[the margin ratio]
The author also asked how the cash-line formulation of an earlier version (Remark~\ref{rem:oldCP} in the
Supplement) could be written with the loss ratio alone. Under the representation,
$1-\LR(X+c)=(\Pi(X)-\rho(X))/(\Pi(X)+c)$, so off the break-even set the map $c\mapsto1/(1-\LR(X+c))$ is affine; it
vanishes at the price point and equals $1$ where the cash line vanishes. The author proposed to call
$\mathrm{MR}:=1-\LR$ the \emph{margin ratio}, the underwriting margin per unit of premium, and to state the
condition through the ratio of the two: since $1/(1-\LR)=1+\LR/\mathrm{MR}$, every cash line off the break-even set
is a unit-slope Möbius map if and only if $c\mapsto\LR(X+c)/\mathrm{MR}(X+c)$ is affine, with slope
$1/(\Pi(X)-\rho(X))$. The author's first version of the identity had the quotient inverted, and the main assistant
corrected it. Both forms degenerate on the break-even set, where the margin ratio vanishes along the whole cash line
except at the price point, so they do not locate the premium of a break-even book. The condition was therefore kept
as a reading of Proposition~\ref{prop:readings}(a), not as an axiom.
\end{discussion}

Monotonicity is imposed on the cash offsets in \axRM{} and not on the values of the ratio, for the following
reason.

\begin{proposition}[No monotone loss ratio]\label{prop:F2}
Let $h,g\in\Hset_1$ with $g>0$ on $(0,1]$, and let $\LR=I_h/I_g$ on $\Xp:=L^\infty_+\setminus\{0\}$.
If $\LR$ is monotone on $\Xp$ (non-decreasing or non-increasing for the almost sure order),
then $h=g$ and $\LR\equiv1$.
\end{proposition}

\begin{proof}
Let $P(A)=p>0$. For $c\ge0$, $\LR(\1_A+c)=\frac{h(p)+c}{g(p)+c}$ is monotone in $c$ with derivative at
$0$ of the sign of $g(p)-h(p)$. Next let $P(E)=q>0$ and $t\ge0$. Then $X_t=1+t\1_E\in\Xp$ is
increasing in $t$, and $I_k(X_t)=1+t\,k(q)$ for $k\in\{h,g\}$, so
$\LR(X_t)=\frac{1+t\,h(q)}{1+t\,g(q)}$ has derivative at $t=0$ of the sign of $h(q)-g(q)$. If $\LR$
is non-decreasing both signs are non-negative, hence $h=g$ on $(0,1]$; if $\LR$ is non-increasing both are
non-positive, and again $h=g$. \qedhere
\end{proof}
